The bivector element of $\Cl(V_\C)$ attached to $a\wedge b$ is $\tfrac12 c(a)c(b) - \tfrac12\ip{a}{b}$ (the product of two vectors is $\ip{a}{b}$ plus half their commutator, \claimref{std-algebra.vector-squares}): its commutator with vectors is the bivector action for any $a, b$, orthogonal or not, as in Proposition \claimref{correspondence.spin-representation}. Take $v = w_1$, $v^* = w_1'$ in an adapted model, and write $A_1 := w_1\wedge(\cdot)$, $B_1 := \iota_{w_1'}$, so $c(v)c(v^*) = 2\,A_1 B_1$. On a basis monomial containing $w_1$, $A_1B_1$ removes and restores the $w_1$-factor: eigenvalue $1$; on a monomial without $w_1$, $B_1$ annihilates it: eigenvalue $0$. So $\tfrac12 c(v)c(v^*) = N$ and $\rho(E) = N - \tfrac12$. The script checks the full $8\times 8$ eigenvalue table. The eigenvalues $\pm\tfrac12$ on the two chirality-plane pairs are the grading weights used in Proposition \claimref{correspondence.local-twistor-transport-ascended-spinor-connection}.
